\par\noindent\textbf{(1)} 频域 Maxwell 方程 k×E=ωB、k×H=−ωD、D=ε·E，消去H得 k×(k×E)+ω²μ₀ε·E=0。对 k=(k sinθ,0,k cosθ)，y分量独立，xz分量形成二维耦合系统；还须满足 k·D=0，而不是假设 k·E=0。
\par\noindent\textbf{(2)} y偏振的寻常波满足 k²=ω²μ₀ε\_o。xz面偏振的非寻常波由行列式给 k\_x²/ε\_e+k\_z²/ε\_o=μ₀ω²。于是 v\_\{p,e\}\textasciicircum{}\{−2\}=μ₀(sin²θ/ε\_e+cos²θ/ε\_o)，其折射率随波法线角度变化。
\par\noindent\textbf{(3)} 在色散曲面 ω(k)=μ₀\textasciicircum{}\{−1/2\}√(k\_x²/ε\_e+k\_z²/ε\_o)，群速度分量为 ∂ω/∂k\_x=k\_x/(μ₀ε\_eω)、∂ω/∂k\_z=k\_z/(μ₀ε\_oω)。因此 v\_\{g,x\}/v\_\{g,z\}=(ε\_o/ε\_e)tanθ，除光轴、垂直轴和 ε\_e=ε\_o 外群速不沿k，能流在晶体内部产生walk-off。
\par\noindent\textbf{(4)} 沿光轴θ=0时两模均为k²=μ₀ε\_oω²，简并；θ=π/2 时非寻常有效介电为 ε\_e，而寻常为 ε\_o。当ε\_e→ε\_o，两个色散面合一且群速平行k，回复各向同性介质。以上推导忽略磁各向异性和吸收。
